% TEMPLATE-MILC-v3.tex - plantilla maestra UNICA de las guias MILC v3.
%
% La usan las tres familias de guias, cada una con su driver:
%   · Grados 6-11  -> scripts/build-guias-g11.py
%   · Bebras       -> scripts/build-guias-bebras.py
%   · Semillero    -> scripts/build-guias-semillero.py
%
% Antes habia tres copias casi identicas (~400 lineas iguales, 6 distintas):
% cada mejora habia que aplicarla tres veces, y en la practica no se hacia
% (el motor de recursos vivia solo en dos de ellas). Lo que varia entre
% familias esta parametrizado en 6 placeholders que cada driver llena
% (se nombran aqui SIN su sintaxis real: el builder reemplaza por texto plano,
% tambien dentro de los comentarios, y un valor multilinea rompe el preambulo):
%
%   HEADER_IZQ        encabezado izquierdo de pagina
%   HEADER_DER        encabezado derecho de pagina
%   PORTADA_ETIQUETA  rotulo sobre el numero grande (GRADO/MOMENTO/MODULO)
%   PORTADA_SERIE     linea de serie de la portada
%   PORTADA_ID        identificador de la guia en la portada
%   PORTADA_CIRCULO   el circulito de grado (°), o vacio si no aplica
%
% El resto de claves las documenta content/guias/_SCHEMA.md. El script valida
% que no queden placeholders sin reemplazar antes de escribir el archivo final.

\documentclass[11pt,letterpaper]{article}
\usepackage[margin=2.0cm]{geometry}
\usepackage[table]{xcolor}
\usepackage{fontspec}
\usepackage[spanish]{babel}
\usepackage{tikz}
% Librerías TikZ para los diagramas de las guías (bloque `recursos:`):
%  · babel  → babel en español vuelve activos caracteres como «>», y TikZ los usa
%             en sus opciones (>=stealth, ->). Sin esta librería, cualquier
%             diagrama con flechas revienta con «\language@active@arg> has an
%             extra }». Debe ir ANTES de que se dibuje nada.
%  · positioning        → sintaxis «below left=2pt and 3pt».
%  · decorations.pathreplacing → llaves ({brace}) para señalar rangos.
\usetikzlibrary{babel,positioning,decorations.pathreplacing}
\usepackage{graphicx}
% Circuitos: el trabajo de electrónica se hace hoy con micro:bit y sus sensores
% integrados (MakeCode, sin cableado), así que no se necesitan esquemáticos.
% Si algún día entran montajes con protoboard, basta descomentar esta línea:
% los diagramas .tex/.tikz declarados en `recursos:` ya se incrustan solos.
% \usepackage{circuitikz}
\usepackage[most]{tcolorbox}
\usepackage{tabularx}
\usepackage{array}
\usepackage{fancyhdr}
\usepackage{titlesec}
\usepackage[document]{ragged2e}  % bandera a la derecha en todo el cuerpo
\usepackage{enumitem}
\usepackage{microtype}

% Helvetica Neue no trae la flecha U+2192, asi que cada «→» escrito literalmente
% en el YAML se compilaba a NADA, en silencio (462 flechas en 61 guias: rutas del
% tipo «paleta → tipografia → lockup» salian rotas). Mapeamos el caracter a su
% version matematica, que si tiene glifo. Asi el autor puede seguir escribiendo
% «→» comodamente en el YAML.
\usepackage{newunicodechar}
\newunicodechar{→}{\ensuremath{\rightarrow}}
% Mismo problema con los simbolos de marca y vinetas: el «✓» de «una palomita ✓»
% salia como cuadrito vacio en el PDF de todas las guias que lo usaban.
\usepackage{pifont}
\newunicodechar{✓}{\ding{51}}
\newunicodechar{✗}{\ding{55}}
\newunicodechar{✔}{\ding{52}}
\newunicodechar{•}{\textbullet}
\newunicodechar{≥}{\ensuremath{\geq}}
\newunicodechar{≤}{\ensuremath{\leq}}

\setmainfont{Helvetica Neue}
\setsansfont{Helvetica Neue}
\setlength{\parindent}{0pt}
\setlength{\parskip}{6pt}
% Sin particion de palabras: en 6.º-7.º una palabra cortada por guion al final
% de linea («Comuni-cacion») frena la lectura mucho mas que una linea corta.
\hyphenpenalty=10000
\exhyphenpenalty=10000
\pretolerance=2000
\tolerance=3000
\setlength{\emergencystretch}{3em}
\linespread{1.10}
\renewcommand{\arraystretch}{1.25}
\setlist{leftmargin=5mm,itemsep=1mm,topsep=1mm}
\newcolumntype{Y}{>{\RaggedRight\arraybackslash}X}

\definecolor{milcMagenta}{HTML}{E6007E}
\definecolor{milcVino}{HTML}{5A0038}
\definecolor{milcTurquesa}{HTML}{0093A5}
\definecolor{milcVerde}{HTML}{58A923}
\definecolor{milcMostaza}{HTML}{E5B400}
\definecolor{milcOcre}{HTML}{B86600}
\definecolor{milcNegro}{HTML}{241816}
\definecolor{milcGris}{HTML}{F7F7F4}
\definecolor{milcLinea}{HTML}{E8E2DD}
\definecolor{milcGradePrimary}{HTML}{84CC16}
\definecolor{milcGradeDark}{HTML}{4D7C0F}
\definecolor{milcGradeSoft}{HTML}{ECFCCB}
\definecolor{milcGradeAccent}{HTML}{CFE7FF}
\definecolor{milcGradeText}{HTML}{FFFFFF}
\color{milcNegro}

\pagestyle{fancy}
\fancyhf{}
\lhead{\small Tecnología e Informática · Grado 7}
\rhead{\small Guía 19 · MILC v3}
\cfoot{\small\thepage}
\renewcommand{\headrulewidth}{0pt}

\newtcolorbox{titlebox}[1]{
  enhanced,colback=#1,colframe=#1,arc=7mm,boxrule=0pt,
  left=7mm,right=7mm,top=3mm,bottom=3mm,
  before skip=8pt,after skip=8pt,
  fontupper=\sffamily\bfseries\Large\color{white}
}

\newtcolorbox{softbox}[3]{
  enhanced,colback=#2,colframe=#1,borderline west={4pt}{0pt}{#1},
  arc=2mm,boxrule=.4pt,left=4mm,right=4mm,top=3mm,bottom=3mm,
  before skip=6pt,after skip=8pt,title={#3},coltitle=white,
  colbacktitle=#1,fonttitle=\sffamily\bfseries,fontupper=\RaggedRight,
  attach boxed title to top left={xshift=3mm,yshift=-2mm},
  boxed title style={arc=2mm, boxrule=0pt}
}

\newtcolorbox{drawbox}[2]{
  enhanced,colback=white,colframe=#1,arc=4mm,boxrule=1pt,
  left=5mm,right=5mm,top=4mm,bottom=4mm,before skip=8pt,after skip=8pt,
  title={#2},coltitle=white,colbacktitle=#1,fonttitle=\sffamily\bfseries,
  fontupper=\RaggedRight,
  attach boxed title to top left={xshift=4mm,yshift=-2mm},
  boxed title style={arc=3mm, boxrule=0pt}
}

\newtcolorbox{infoband}[1]{
  enhanced,colback=#1!8,colframe=#1!50,arc=2mm,boxrule=.5pt,
  left=4mm,right=4mm,top=2mm,bottom=2mm,before skip=4pt,after skip=4pt,
  fontupper=\small\RaggedRight
}

% Puente narrativo entre secciones (contrato editorial Conéctate).
% Da continuidad: "Acabas de X. Ahora vas a Y porque Z."
% Visualmente: banda izquierda en color del grado, texto en itálica pequeña.
\newtcolorbox{puentebox}{
  enhanced,colback=milcGradeSoft,colframe=milcGradeDark,
  borderline west={3pt}{0pt}{milcGradeDark},
  arc=0mm,boxrule=0pt,left=5mm,right=4mm,top=2.5mm,bottom=2.5mm,
  before skip=4pt,after skip=8pt,
  fontupper=\itshape\small\color{milcNegro}\RaggedRight
}
% La flecha va en modo matematico a proposito: Helvetica Neue NO tiene el glifo
% U+2192, asi que \textrightarrow se compilaba a NADA («Missing character: There
% is no → in font Helvetica Neue») y el puente salia sin flecha. En modo
% matematico la flecha viene de la fuente matematica y si se dibuja.
\newcommand{\puente}[1]{%
  \begin{puentebox}%
    \textcolor{milcGradeDark}{$\rightarrow$}\hspace{2mm}#1%
  \end{puentebox}%
}

\newcommand{\linead}[1]{\noindent\color{milcLinea}\rule{#1}{0.5pt}\color{milcNegro}}
\newcommand{\respuesta}[1]{\par\vspace{2mm}\foreach \n in {1,...,#1}{\noindent\color{milcLinea}\rule{\linewidth}{0.5pt}\par\vspace{5mm}}\color{milcNegro}}
\newcommand{\checkbox}{\textcolor{milcGradeDark}{\fbox{\phantom{x}}}\hspace{5pt}}

\newcommand{\phasecard}[4]{
\begin{tcolorbox}[enhanced,colback=#3,colframe=#2,arc=3mm,boxrule=.5pt,height=3.05cm,valign=center,halign=center,left=1.5mm,right=1.5mm,top=2mm,bottom=2mm]
{\sffamily\bfseries\fontsize{22}{24}\selectfont\textcolor{#2}{#1}}\par
{\sffamily\bfseries\small #4}
\end{tcolorbox}
}

% ── Piezas del rediseno 2026 (legibilidad 6.º-11.º) ───────────────────
% Banda de estacion: reemplaza el titlebox macizo. Titulo a la izquierda,
% dato practico (tiempo/modalidad) a la derecha, en una sola linea.
\newcommand{\milcestacion}[2]{%
\begin{tcolorbox}[enhanced,colback=#1,colframe=#1,arc=2mm,boxrule=0pt,
  left=5mm,right=5mm,top=2.6mm,bottom=2.6mm,before skip=11pt,after skip=7pt]
{\sffamily\bfseries\fontsize{15}{18}\selectfont\color{white}#2}
\end{tcolorbox}}

% Tarjeta de paso numerada: sustituye las tablas «Pilar / Aplicacion», que
% eran formato de consulta docente, no de estudiante. El numero va en un
% circulo sobre el borde para que la secuencia se lea de un vistazo.
\newcommand{\milcpaso}[3]{%
\begin{tcolorbox}[enhanced,colback=white,colframe=milcLinea,arc=1.5mm,boxrule=.9pt,
  left=14mm,right=4mm,top=3mm,bottom=3mm,before skip=4pt,after skip=4pt,
  fontupper=\RaggedRight,
  overlay={\node[anchor=north west,circle,fill=#2,text=white,inner sep=0pt,
    minimum size=7.4mm,font=\sffamily\bfseries\small]
    at ([xshift=3.6mm,yshift=-3.2mm]frame.north west) {#1};}]
#3
\end{tcolorbox}}

% Casilla marcable de tamano util con lapiz (la anterior era \fbox{\phantom{x}},
% de alto variable segun el texto que la seguia).
\newcommand{\milcmarca}{\framebox[3.8mm]{\rule{0pt}{3.4mm}}}

% Fila marcable de lista suelta (checks de la estacion 1).
\newcommand{\milccheckrow}[1]{%
\par\vspace{1.8mm}\noindent
\makebox[6.4mm][l]{\raisebox{-.55mm}{\textcolor{milcNegro}{\milcmarca}}}%
\parbox[t]{\dimexpr\linewidth-6.4mm\relax}{\RaggedRight #1}\par}

% Celda marcable dentro de la rubrica (tabla de autoevaluacion). Misma
% sangria francesa, pero pensada para vivir dentro de una columna X.
\newcommand{\milcopcion}[1]{%
\makebox[5.2mm][l]{\raisebox{-.55mm}{\textcolor{milcNegro}{\milcmarca}}}%
\parbox[t]{\dimexpr\linewidth-5.2mm\relax}{\RaggedRight #1}}

% ── Recursos visuales de la guía ──────────────────────────────────────
% Declarados en el bloque `recursos:` del YAML y emitidos por el builder.
% \guiaFigura   → imagen rasterizada o PDF (foto, carta celeste, montaje).
% \guiaDiagrama → fuente TikZ/circuitikz (.tex) incrustada como vector:
%                 es la vía para circuitos y esquemáticos editables.
% #1 = ruta relativa al .tex · #2 = pie (puede ir vacío)
\newcommand{\guiaFigura}[2]{%
\begin{center}
\includegraphics[width=0.88\linewidth,height=0.38\textheight,keepaspectratio]{#1}\\[1.5mm]
{\footnotesize\itshape\textcolor{milcNegro!75}{#2}}
\end{center}
}

\newcommand{\guiaDiagrama}[2]{%
\begin{center}
\input{#1}\\[1.5mm]
{\footnotesize\itshape\textcolor{milcNegro!75}{#2}}
\end{center}
}

\begin{document}
\thispagestyle{empty}

% ── Portada ───────────────────────────────────────────────────────────
% Antes: un rectangulo de color a pagina completa. Se veia bien en pantalla,
% pero en la fotocopiadora del colegio se comia el toner y salia gris sucio.
% Ahora la banda ocupa el tercio superior y el resto va en blanco.
\begin{tikzpicture}[remember picture,overlay]
  \path[fill=milcGradePrimary]
    (current page.north west) rectangle ([yshift=-10.6cm]current page.north east);

  \node[anchor=north west,text=milcGradeText] at ([xshift=2.0cm,yshift=-1.55cm]current page.north west)
    {\sffamily\bfseries\fontsize{12}{14}\selectfont GRADO};
  \node[anchor=north west,text=milcGradeText,opacity=.85] at ([xshift=2.0cm,yshift=-2.35cm]current page.north west)
    {\sffamily\bfseries\fontsize{8.5}{10}\selectfont SERIE GUÍAS · TECNOLOGÍA E INFORMÁTICA};
  \node[anchor=north east,text=milcGradeText,opacity=.85,align=right] at ([xshift=-2.0cm,yshift=-1.55cm]current page.north east)
    {\sffamily\bfseries\fontsize{9}{12}\selectfont I.E. Sor María Juliana\\Cartago, Valle del Cauca};

  \node[anchor=west,text=milcGradeText] (gradonum) at ([xshift=1.85cm,yshift=-7.1cm]current page.north west)
    {\sffamily\bfseries\fontsize{112}{112}\selectfont 7};
  % El circulito de grado (°) solo aplica a las guias de grado («11°»); en
  % Bebras y Semillero el numero grande es un momento/modulo, no un grado.
  % Se posiciona relativo al nodo `gradonum`, no en coordenadas absolutas.
  \draw[milcGradeText,line width=1.6pt]
    ([xshift=2.0mm,yshift=-4.5mm]gradonum.north east) circle (.15cm);

  \node[anchor=west,text=milcGradeText,text width=11.4cm,align=left] at ([xshift=1.5cm,yshift=0.5cm]gradonum.east)
    {
     {\sffamily\bfseries\fontsize{9}{11}\selectfont\MakeUppercase{Período 2 · Algoritmia y pensamiento computacional}}\\[3mm]
     {\sffamily\bfseries\fontsize{21}{25}\selectfont\setlength{\baselineskip}{27pt}Scratch ·\\[4pt]narrativa\\[4pt]interactiva ---\\[4pt]el cuento\\[4pt]que escoge\\[4pt]el lector}\\[3.5mm]
     {\sffamily\bfseries\fontsize{15}{18}\selectfont Guía 19-7-TIC}
    };
\end{tikzpicture}

\vspace*{9.4cm}
\vfill

{\sffamily\bfseries\fontsize{8}{10}\selectfont\textcolor{milcGradeDark}{LO QUE TE LLEVAS HOY}}

\begin{tcolorbox}[enhanced,colback=white,colframe=milcNegro,arc=2mm,boxrule=1.6pt,
  left=5mm,right=5mm,top=4mm,bottom=4mm,before skip=3pt,after skip=12pt,
  fontupper=\RaggedRight\fontsize{11}{15.5}\selectfont]
Una \textbf{narrativa interactiva en Scratch} de mínimo \textbf{3 escenas}. La historia: \textbf{(a) inicia con un sprite que cuenta el inicio}; \textbf{(b) presenta al lector 2 opciones (preguntar con bloque de Sensores)}; \textbf{(c) según la opción, va a una escena u otra} (usando condicional ``si...entonces''); \textbf{(d) cada escena tiene fondo distinto, sprite cambia disfraz, sonidos opcionales}. Más \textbf{el guion previo escrito en cuaderno} (esquema de la historia + 2 caminos) + \textbf{captura del proyecto guardado}.
\end{tcolorbox}

\begin{tcolorbox}[enhanced,colback=milcGris,colframe=milcGris,arc=2mm,boxrule=0pt,
  left=5mm,right=5mm,top=3.5mm,bottom=3.5mm,after skip=12pt]
{\sffamily\bfseries\fontsize{8}{10}\selectfont\textcolor{milcNegro!55}{ESTA GUÍA ES DE}}\\[3mm]
\begin{tabularx}{\linewidth}{X p{3.2cm} p{3.2cm}}
\linead{\linewidth} & \linead{3.0cm} & \linead{3.0cm}\\[-1mm]
{\fontsize{8.5}{10}\selectfont\textcolor{milcNegro!55}{Nombre completo}} &
{\fontsize{8.5}{10}\selectfont\textcolor{milcNegro!55}{Grupo}} &
{\fontsize{8.5}{10}\selectfont\textcolor{milcNegro!55}{Fecha}}\\
\end{tabularx}
\end{tcolorbox}

\vfill

\noindent\textcolor{milcLinea}{\rule{\linewidth}{1pt}}\\[2mm]
{\sffamily\bfseries\fontsize{9.5}{12}\selectfont Autor: Dr. Álvaro Cárdenas Orozco}\\[1mm]
{\sffamily\fontsize{8.5}{11}\selectfont\textcolor{milcNegro!55}{Metodología MILC · apertura ancestral · pensamiento computacional · triángulo Dussel–estoicismo–Floridi}}

\newpage

\section*{Apertura: Scratch · narrativa interactiva --- el cuento que el lector escoge}

\begin{softbox}{milcGradeDark}{milcGradeSoft}{Saber ancestral}
\textbf{Doña Mercedes la maestra rural de Cartago narraba cuentos donde los niños decidían.} En la escuelita de la vereda La Plata, los miércoles eran días de \textbf{cuento abierto}. Doña Mercedes empezaba: \emph{``Iba un pescador caminando por el río Cauca cuando vio dos caminos''}. Luego preguntaba: \emph{``¿Va por la izquierda o por la derecha?''}. Los niños gritaban \emph{``¡Izquierda!''}. Doña Mercedes seguía: \emph{``Por la izquierda encontró un pájaro herido. ¿Lo ayuda o sigue?''}. Los niños decidían otra vez. \textbf{Cada cuento era distinto porque el grupo decidía qué pasaba}. Doña Mercedes no improvisaba sin orden: \emph{tenía un mapa mental con 4 o 5 caminos posibles, cada uno con su propio final}. Los niños se sentían \emph{coautores} del cuento. Esa pedagogía interactiva existía antes de los videojuegos y antes de Scratch. \textbf{Hoy haces lo mismo: programas un cuento donde el lector decide los caminos}.
\end{softbox}

\begin{softbox}{milcTurquesa}{milcGris}{Saber contemporáneo}
Una \textbf{narrativa interactiva} (o \emph{historia ramificada}) es un cuento donde el lector toma decisiones que cambian el rumbo. En lugar de un único final, hay varios. En lugar de leer pasivo, el lector \textbf{participa}. Esta forma narrativa existe en libros (\emph{``Escoge tu propia aventura''} de los años 80), en videojuegos (Detroit Become Human, Life is Strange), en películas (Black Mirror Bandersnatch). \textbf{En Scratch puedes hacer la tuya} usando: \emph{(a) bloque Sensores ``preguntar [pregunta] y esperar''} (recibe respuesta del lector); \emph{(b) bloque Sensores ``respuesta''} (guarda lo que el lector escribió); \emph{(c) bloque Control ``si... entonces... si no''} (decide rumbo según respuesta); \emph{(d) bloques Apariencia ``cambiar fondo'' y ``cambiar disfraz''} (cambian escena). Hoy combinas variables (S5), condicionales (S6) y Scratch (S8) en una narrativa que tu equipo va a vivir.
\end{softbox}

\begin{softbox}{milcMostaza}{milcGris}{Pregunta-puente}
\textit{Si pudieras programar una historia donde el lector escoge qué hace el protagonista, ¿\textbf{qué tema escogerías}? ¿Un héroe wayuu? ¿Una niña perdida en el campo? ¿Un astronauta? ¿Un cuento de tu vida?}
\end{softbox}

\begin{softbox}{milcVerde}{milcGris}{Saber y saber hacer}
Al terminar la clase: (1) podrás \textbf{identificar} la estructura de una narrativa ramificada; (2) sabrás \textbf{aplicar} los bloques Sensores y Apariencia de Scratch; (3) podrás \textbf{evaluar} la coherencia de una historia interactiva; (4) habrás \textbf{creado} tu primera narrativa interactiva.
\end{softbox}



\small
\begin{tabularx}{\textwidth}{>{\bfseries}p{3.0cm}Y}
\rowcolor{milcGradePrimary!12}Autor & Dr. Álvaro Cárdenas Orozco\\
DBA & Construye algoritmos y diagramas de flujo para resolver problemas paso a paso (MEN, T\&I 7°)\\
Referentes & Saber ancestral del tejedor · Pensamiento computacional · Scratch\\
Producto final & Una \textbf{narrativa interactiva en Scratch} de mínimo \textbf{3 escenas}. La historia: \textbf{(a) inicia con un sprite que cuenta el inicio}; \textbf{(b) presenta al lector 2 opciones (preguntar con bloque de Sensores)}; \textbf{(c) según la opción, va a una escena u otra} (usando condicional ``si...entonces''); \textbf{(d) cada escena tiene fondo distinto, sprite cambia disfraz, sonidos opcionales}. Más \textbf{el guion previo escrito en cuaderno} (esquema de la historia + 2 caminos) + \textbf{captura del proyecto guardado}.\\
\end{tabularx}
\normalsize

\puente{Hoy haces 4 cosas: planeas tu historia con caminos, aprendes los bloques nuevos, la programas en Scratch, la pruebas con un compañero.}

\milcestacion{milcGradeDark}{Ruta MILC: así vamos a aprender}
\begin{tabularx}{\textwidth}{>{\centering\arraybackslash}X>{\centering\arraybackslash}X>{\centering\arraybackslash}X>{\centering\arraybackslash}X}
\phasecard{1}{milcVerde}{milcGris}{Escucha\\[-1mm]\small Observo, escucho y pregunto.} &
\phasecard{2}{milcTurquesa}{milcGris}{Entender\\[-1mm]\small Sistematizo y ordeno.} &
\phasecard{3}{milcMagenta}{milcGris}{Hacer\\[-1mm]\small Construyo y aplico.} &
\phasecard{4}{milcVino}{milcGris}{Cierre\\[-1mm]\small Evalúo y reflexiono.}
\end{tabularx}


\puente{Empezamos diseñando el árbol de tu historia en papel.}

\milcestacion{milcVerde}{Estación 1 de 3 · Escucha}

\begin{softbox}{milcVerde}{milcGris}{Escena para escuchar}
\textbf{Actividad 1 · ANALIZA --- Diseña el árbol de tu historia} (12 min · individual). Antes de tocar Scratch, planea la historia en el cuaderno. En una hoja en blanco, dibuja el \textbf{árbol de tu historia} con esta estructura: \emph{(a) Escena inicial:} situación + protagonista. \emph{(b) Pregunta al lector con 2 opciones.} \emph{(c) Escena A si elige opción 1.} \emph{(d) Escena B si elige opción 2.} Cada rama termina con su propio final. Sugerencias de temas: \emph{``Una niña en el campo encuentra dos caminos''}; \emph{``Un astronauta llega a un planeta nuevo''}; \emph{``Un detective tiene 2 pistas para escoger''}; \emph{``Un héroe wayuu enfrenta una decisión''}. Hazlo simple: 3 escenas, 2 caminos, 2 finales.
\end{softbox}

{\sffamily\bfseries\fontsize{8}{10}\selectfont\textcolor{milcNegro!55}{MARCA CUANDO LO HAYAS HECHO}}
\milccheckrow{Dibujé el árbol de la historia.}
\milccheckrow{Tiene escena inicial + pregunta + 2 ramas + 2 finales.}
\milccheckrow{Cada rama tiene su propio fondo y acción.}

\begin{infoband}{milcVerde}


\textbf{Lo que va al cuaderno (Actividad 1 · ANALIZA).} Título: «Actividad 1 --- Árbol de mi historia interactiva». \textbf{Formato:} dibujo del árbol con 3+ escenas y 2 caminos. \textbf{Extensión:} 1/2 página. \textbf{Sabes que terminaste cuando} tu árbol es claro y podrías programarlo sin reinventarlo.
\end{infoband}

\puente{Después aprendes los bloques nuevos: preguntar, respuesta, cambiar fondo.}

\milcestacion{milcTurquesa}{Estación 2 de 3 · Entender}

\begin{softbox}{milcTurquesa}{milcGris}{Lo que vas a entender}
Tu árbol es el guion. Ahora aprendes los \textbf{bloques nuevos de Scratch} que necesitas para implementarlo.
\end{softbox}

\milcpaso{1}{milcTurquesa}{\textbf{Bloque Sensores --- ``Preguntar y esperar''.} Color turquesa. Estructura: \emph{``Preguntar [tu pregunta] y esperar''}. Cuando se ejecuta este bloque, aparece una caja de texto en la parte inferior del escenario donde el lector escribe su respuesta. El programa se pausa hasta que el lector escriba y pulse Enter. Es como una conversación con el usuario.}
\milcpaso{2}{milcTurquesa}{\textbf{Bloque Sensores --- ``Respuesta''.} También turquesa. Es una variable especial que guarda lo último que el lector escribió en una pregunta. Si preguntaste \emph{``¿Vas a la izquierda o derecha?''} y el lector escribió \emph{``izquierda''}, entonces \texttt{respuesta = izquierda}. \textbf{Cómo se usa en condicional:} \emph{Si respuesta = ``izquierda'' entonces ir a escena A; si no, ir a escena B}.}
\milcpaso{3}{milcTurquesa}{\textbf{Bloques Apariencia --- ``Cambiar fondo'' y ``Cambiar disfraz''.} Color morado. Para que tu historia tenga diferentes escenas, necesitas cambiar el \textbf{fondo} del escenario (\emph{bosque, ciudad, espacio}). El bloque \emph{``Cambiar fondo a [nombre]''} lo hace. También para que el sprite cambie de aspecto: \emph{``Cambiar disfraz a [nombre]''} (el sprite puede tener varios disfraces: \emph{feliz, triste, asustado}).}
\milcpaso{4}{milcTurquesa}{\textbf{Estructura típica de una escena.} Cada escena en tu programa sigue este patrón: \emph{(1) Cambiar fondo}, \emph{(2) Cambiar disfraz del sprite}, \emph{(3) Decir texto narrativo}, \emph{(4) Preguntar (si aplica)}, \emph{(5) Decisión con condicional (si aplica)}. Cuando ya conoces este patrón, armar cada escena es \emph{repetir la estructura} con datos distintos. Eso es \emph{patrones} de la S3 aplicados a Scratch.}

\begin{drawbox}{milcTurquesa}{Bloques nuevos + estructura de escena}
\textbf{Sensores (turquesa):} preguntar y esperar · respuesta. \\[1mm]
\textbf{Apariencia (morado):} cambiar fondo · cambiar disfraz. \\[1mm]
\textbf{Estructura de escena:} cambiar fondo · cambiar disfraz · decir narrativa · preguntar · condicional.
\end{drawbox}

\begin{softbox}{milcMostaza}{milcGris}{Errores comunes y su corrección}
\textbf{Errores frecuentes en narrativas interactivas.} (1) \emph{Olvidar el bloque ``preguntar y esperar''} antes de leer la respuesta --- el programa no espera al lector. (2) \emph{Condicional con respuesta exacta} (\emph{``si respuesta = Izquierda''}) pero el lector escribió \emph{``izquierda''} con minúscula --- no coincide. Solución: usar \emph{``mayúscula a minúscula''} o aceptar variaciones. (3) \emph{Olvidar cambiar fondo entre escenas} --- todas las escenas se ven igual. (4) \emph{Solo 1 final feliz} --- las narrativas interactivas son ricas cuando cada decisión tiene consecuencia real. (5) \emph{Demasiado texto para leer} --- los lectores se aburren; mejor narrativa con elecciones rápidas que con monólogos largos.
\end{softbox}

\begin{infoband}{milcTurquesa}


\textbf{Lo que va al cuaderno (Actividad 2 · EXPLICA).} Título: «Actividad 2 --- Bloques nuevos para narrativa». \textbf{Formato:} lista de bloques nuevos con su función + estructura de escena en plantilla. \textbf{Extensión:} 10 líneas. \textbf{Sabes que terminaste cuando} podrías diseñar cualquier escena de Scratch siguiendo el patrón.
\end{infoband}

\puente{Con todo claro, programas la historia en Scratch.}

\milcestacion{milcMagenta}{Estación 3 de 3 · Hacer}

\begin{softbox}{milcMagenta}{milcGris}{Cómo vas a construir tu producto}
\textbf{Actividad 3 · CREA --- Tu narrativa interactiva en Scratch} (35 min · individual). Aplica tu árbol de la Actividad 1 a Scratch. Es un trabajo más ambicioso que el de la S8 pero usando los mismos componentes + los nuevos bloques de hoy.
\end{softbox}

\milcpaso{1}{milcMagenta}{\textbf{Paso 1 --- Abre Scratch y guarda con buen nombre.} En \texttt{scratch.mit.edu} crea proyecto nuevo. Si tienes cuenta, renómbralo: \texttt{2026-mi-historia-interactiva-[tu-nombre]}. Guarda. \textbf{Antes de programar, ten tu árbol del cuaderno a la vista}: es tu guion.}
\milcpaso{2}{milcMagenta}{\textbf{Paso 2 --- Selecciona o sube fondos para las 3 escenas.} Botón \emph{``Elegir un fondo''} (abajo a la derecha). Escoge 3 fondos distintos (\emph{bosque, río, ciudad}) o sube los tuyos. Los fondos representan las 3 escenas de tu historia. Verifica que aparecen los 3 en la lista de fondos.}
\milcpaso{3}{milcMagenta}{\textbf{Paso 3 --- Programa la Escena 1 (inicial).} En el sprite (gato u otro), arma: \emph{(a) Al hacer clic en bandera verde [Eventos].} \emph{(b) Cambiar fondo a [fondo 1] [Apariencia].} \emph{(c) Decir [primer texto narrativo] por 3 segundos [Apariencia].} \emph{(d) Decir [segundo texto narrativo] por 3 segundos.} \emph{(e) Preguntar [pregunta con 2 opciones] y esperar [Sensores]}. Por ejemplo: \emph{``¿Vas a la IZQUIERDA o a la DERECHA? Escribe izquierda o derecha''}.}
\milcpaso{4}{milcMagenta}{\textbf{Paso 4 --- Programa el condicional con la respuesta.} Justo después del bloque ``preguntar y esperar'', agrega: \emph{Si [respuesta = ``izquierda''] entonces [llamar a Escena A] si no [llamar a Escena B]}. Aquí ``llamar a Escena A'' significa: cambiar fondo a [fondo 2], cambiar disfraz, decir narrativa de la escena A.}
\milcpaso{5}{milcMagenta}{\textbf{Paso 5 --- Programa Escena A y Escena B.} Cada una tiene su propio conjunto de bloques: cambiar fondo → cambiar disfraz → decir texto narrativo → decir final. Si quieres, agrega un sonido al final con bloque ``Tocar sonido''. Cada escena debe terminar la historia (no dejar el lector a medias).}
\milcpaso{6}{milcMagenta}{\textbf{Paso 6 --- Prueba con compañero + captura + reflexión.} Pídele a un compañero que juegue tu historia (haga clic en bandera verde, escriba sus respuestas). Verifica que los 2 caminos funcionan. Si algo falla, ajusta. En el cuaderno: captura del proyecto + \emph{``El camino que más me gustó fue \_\_\_. La parte que más me costó programar fue \_\_\_''}. Firma con nombre y fecha.}

\begin{drawbox}{milcMagenta}{Antes de cerrar la S9}
\checkbox Tengo el árbol de la historia dibujado en el cuaderno. \\[1mm]
\checkbox Mi proyecto Scratch tiene 3 fondos distintos. \\[1mm]
\checkbox La Escena 1 tiene narrativa + pregunta al lector. \\[1mm]
\checkbox El condicional dirige a Escena A o Escena B según la respuesta. \\[1mm]
\checkbox Cada escena tiene su propio fondo, narrativa y final. \\[1mm]
\checkbox Un compañero jugó la historia y los 2 caminos funcionan.
\end{drawbox}

\begin{softbox}{milcMagenta}{milcGris}{Plantilla de trabajo}
\textbf{Estructura del programa.} \textit{Escena 1:} cambiar fondo + decir + preguntar. \textit{Condicional:} si respuesta = opción 1 → Escena A; si no → Escena B. \textit{Escena A:} cambiar fondo + cambiar disfraz + decir + final. \textit{Escena B:} igual estructura. \textit{Cuaderno:} árbol + captura + reflexión + firma.
\end{softbox}

\begin{infoband}{milcMagenta}


\textbf{Lo que va al cuaderno (Actividad 3 · CREA).} Título: «Actividad 3 --- Mi narrativa interactiva». \textbf{Formato:} árbol de historia + captura del proyecto + reflexión. \textbf{Extensión:} 1 página. \textbf{Sabes que terminaste cuando} tu historia se puede jugar de inicio a fin por los 2 caminos.
\end{infoband}

\puente{El producto es la historia jugable + guion en cuaderno.}

\milcestacion{milcMostaza}{Producto de la sesión}
\textbf{Narrativa interactiva en Scratch · 3 escenas · 2 caminos · jugada por compañero}.
\begin{softbox}{milcMostaza}{milcGris}{Criterios de calidad}
(1) Hay 3 escenas con fondos distintos. (2) La escena inicial tiene pregunta al lector. (3) El condicional dirige correctamente según la respuesta. (4) Cada camino llega a un final propio. (5) Un compañero jugó la historia exitosamente por los 2 caminos.
\end{softbox}

\puente{Antes de salir, un compañero juega tu historia y verifica que los 2 caminos funcionan.}

\milcestacion{milcVino}{Evaluación liberadora · cinco dimensiones}

\begin{softbox}{milcVino}{milcGris}{Sentido de la evaluación}
Evaluar no es solo revisar una nota. Es reconocer qué aprendí, qué cambió en mí, qué emociones manejé, cómo me relaciono con la ciudad y la comunidad, y qué le dejaré a quien viene detrás.
\end{softbox}

\textbf{Mi reflexión liberadora en cinco dimensiones.} Completa cada frase iniciada:

\small
\begin{tabularx}{\textwidth}{>{\bfseries\RaggedRight\arraybackslash}p{3.4cm}Y>{\RaggedRight\arraybackslash}p{4.6cm}}
\rowcolor{milcVino}\color{white}Dimensión & \color{white}\bfseries Mi respuesta (frame starter) & \color{white}\bfseries Algo que descubrí\\
1. Personal & Lo que más cambió en mí fue \linead{6.5cm} & \linead{4.2cm}\\
2. Emocional & La emoción más fuerte fue \linead{2.5cm} y la regulé \linead{3cm} & \linead{4.2cm}\\
3. Ciudadana & En democracia esto importa porque \linead{6cm} & \linead{4.2cm}\\
4. Local (Cartago) & En mi barrio sería útil para \linead{6.5cm} & \linead{4.2cm}\\
5. Intergeneracional & A mi abuela/abuelo le diría \linead{6.5cm} & \linead{4.2cm}\\
\end{tabularx}
\normalsize

\begin{infoband}{milcVino}
\textbf{Para la próxima sustentación.} Vuelve a esta tabla en seis meses. Verás que las respuestas cambian — no porque lo aprendido cambie, sino porque tú cambias. La evaluación liberadora es una conversación contigo a través del tiempo.
\end{infoband}


\puente{Cerramos con 3 voces sobre por qué dar al lector el poder de decidir es acto creativo.}

\milcestacion{milcGradeDark}{Triángulo de pensamiento · cierre filosófico}

\begin{softbox}{milcVino}{milcGris}{Dussel · lente del nosotros}
\textit{Dar al otro el poder de decidir es el acto más generoso. Tu historia interactiva hace eso con cada lector.}

{\footnotesize\itshape\color{milcNegro!70}--- formulación de esta guía, al modo de Enrique Dussel. No es una cita textual.}

Cuando \textbf{programas una narrativa donde el lector escoge}, no estás \emph{controlando} la experiencia: \emph{la entregas}. Cada lector vivirá la historia que sus elecciones construyen. Algunos llegarán al final feliz, otros al triste. \emph{Esa generosidad creativa --- soltar el control para que otros participen --- es virtud profunda}. Doña Mercedes la sabía con sus cuentos en la vereda. Tú la practicas en Scratch.

\textbf{Cuando creo algo, ¿lo hago para controlar la experiencia del otro o para invitarlo a participar?}
\end{softbox}
\linead{\linewidth}\\[1mm]
\linead{\linewidth}

\begin{softbox}{milcMostaza}{milcGris}{Estoicismo · Marco Aurelio · cuidado interior}
\textit{Cada camino de tu historia debe ser coherente con su propia lógica. No abandones al lector a mitad de un rumbo.}

{\footnotesize\itshape\color{milcNegro!70}--- formulación de esta guía, al modo de Marco Aurelio. No es una cita textual.}

Una narrativa interactiva mal hecha \textbf{abandona al lector}: una rama termina abruptamente, otra contradice la anterior, una decisión no tiene consecuencia. \emph{El estoico construye coherencia en cada rama}: si el lector eligió ir a la izquierda, ahí pasa algo real y la rama tiene su propio final completo. \emph{Esa disciplina narrativa es respeto al lector que confió en tu historia}.

\textbf{Mi historia, ¿trata con respeto las decisiones del lector o lo abandona en alguna rama?}
\end{softbox}
\linead{\linewidth}\\[1mm]
\linead{\linewidth}

\begin{softbox}{milcTurquesa}{milcGris}{Floridi · lente de la infoesfera}
\textit{Las narrativas interactivas son el futuro del cuento. El lector ya no consume: co-crea. Esa es revolución cultural en marcha.}

{\footnotesize\itshape\color{milcNegro!70}--- formulación de esta guía, al modo de Luciano Floridi. No es una cita textual.}

Hace 30 años, los cuentos solo se leían \emph{lineal}: empezabas, leías, terminabas. Hoy las narrativas digitales permiten al lector \textbf{co-crear} la experiencia. \emph{Tú estás del lado correcto de esta revolución}: en grado 7 ya programas narrativas interactivas. En 30 años más, escribir sin interactividad será como hoy escribir a máquina de escribir. \emph{Aprendiste la habilidad temprano}.

\textbf{¿Estoy aprendiendo las formas narrativas del futuro o me quedo en las del pasado?}
\end{softbox}
\linead{\linewidth}\\[1mm]
\linead{\linewidth}

\begin{infoband}{milcNegro}
\textbf{Nota para el docente.} Las tres formulaciones son del autor de la guía, no citas textuales de los pensadores: recogen su lente para pensar el tema, no sus palabras. Si vas a citarlos en clase, remite a la obra original.
\end{infoband}

\milcestacion{milcVino}{Mi autoevaluación · marca una casilla por fila}

{\small
\setlength{\extrarowheight}{2.2pt}
\begin{tabularx}{\textwidth}{>{\RaggedRight\arraybackslash\bfseries}p{3.0cm}YYY}
\rowcolor{milcVino}\color{white}\bfseries Criterio & \color{white}\bfseries Lo logré & \color{white}\bfseries En proceso & \color{white}\bfseries Necesito apoyo\\[.6mm]
Comprensión del tema & \milcopcion{Explico las ideas con mis palabras.} & \milcopcion{Reconozco las ideas, falta precisión.} & \milcopcion{Necesito repasar lo básico.}\\[.8mm]
\rowcolor{milcGris}Pensamiento computacional & \milcopcion{Usé los pilares al armar mi producto.} & \milcopcion{Usé algunos pilares.} & \milcopcion{Necesito releer la estación 2.}\\[.8mm]
Aplicación y producto & \milcopcion{Mi producto cumple los criterios.} & \milcopcion{Está armado, con ajustes pendientes.} & \milcopcion{Aún está en construcción.}\\[.8mm]
\rowcolor{milcGris}Reflexión 5 dimensiones & \milcopcion{Respondí cada una con honestidad.} & \milcopcion{Respondí algunas con detalle.} & \milcopcion{Mis respuestas son muy generales.}\\[.8mm]
Triángulo de pensamiento & \milcopcion{Conecté las 3 voces con mi trabajo.} & \milcopcion{Conecté al menos una.} & \milcopcion{No las relaciono con mi proceso.}\\[.8mm]
\end{tabularx}}

\vspace{3mm}
\textbf{Hoy entendí que:}\\[1mm]
\linead{\linewidth}\\[1mm]
\linead{\linewidth}

\textbf{Mi compromiso para la próxima sesión es:}\\[1mm]
\linead{\linewidth}\\[1mm]
\linead{\linewidth}

\begin{softbox}{milcMostaza}{milcGris}{Cita de despedida · Marco Aurelio}
\textit{``El obstáculo en el camino se vuelve el camino. Lo que estorba al camino se vuelve camino.''}\\[1mm]
Cada error de hoy, bien observado, se convierte en pieza del camino que pisarás mañana.
\end{softbox}

\small
\begin{tabularx}{\textwidth}{>{\bfseries}p{3.6cm}Y}
\rowcolor{milcGradeDark!12}Nombre completo: & \linead{10cm}\\
Fecha: & \linead{4cm} \quad Curso/grupo: \linead{4cm}\\
Mi firma: & \linead{9cm}\\
\end{tabularx}
\normalsize

\begin{infoband}{milcGradeDark}
\textbf{Para la próxima sesión.} Esta semana, mejora tu narrativa: agrega una tercera opción, un sonido, otra escena. Comparte el enlace con familiares para que la jueguen. La próxima clase es \textbf{la cosecha del periodo}: programa completo en Scratch que integra TODO lo aprendido en P2.
\end{infoband}

\end{document}
